\scanpage{56}
\mspage{12}
\fbox{\deleted{\unsure{A)|4)}}} \inserted{(B)}
\uline{Opération} \uline{de} \deleted{\uline{affaiblissement \unsure{et|ou}} de} \inserted{\uline{sur}}\uline{marquage}

\deleted{On peut le définir \inserted{en termes de} par la donnée, d'une \inserted{\unsure{composante connexe de}} pour chaque arête du contour $R$ \unsure{$R\ne\phi$} $\widetilde\Sigma\setminus R$ Soit $\unsure{\widetilde\Sigma_*}\;A_0^+\unsure{\subset|\in}\unsure{A_0|A_*}$ l'ens}
\deleted{\unsure{des}}
Soit $A^+$ l'ens des arêtes \inserted{(libres ou liées)} ``effectivement marquées'', \inserted{du $\Gamma$ bip. telles que $r(a)>0$},
\deleted{et $B$ l'ens des},
correspondant à des $C\in\pi_0(\Sigma)$ telles que $C\cap R\ne\phi$, soit $B^+$ l'ens des comp. connexes \unsure{des|de}
\[
 \Sigma^+\setminus R\cap\Sigma^+
\]
(lesquelles sont des int. ouverts), un surmarquage est défini par la donnée d'une application
\[
 B^+\xrightarrow{r'}\mathbb N,
\]
et peut se décrire \deleted{par} \inserted{comme} l'opération qui sur chaque
\uline{arête} \inserted{\unsure{(ouverte)|\deleted{(ouverte)}} \deleted{qui}}
de $\Sigma$, correspondant à un $b\in B^+$, $r'(b)$ nouveaux sommets de \deleted{\unsure{marquage}} \inserted{rigidification}.
Cette fois-ci, \deleted{\unsure{l'opé}} il y a des choix arbitraires --- néanmoins l'opération donne un objet de $S\mathcal TD^+$, défini à isom unique près \inserted{soit $(X,\Sigma,R_1)$} --- muni bien entendu d'un ens marqué $R'\subset R_1$, tel que l'opération inverse de gommage de $R'$ redonne l'objet de départ $(X,\Sigma,R)$.

\margin{\deleted{arêtes} cercles \deleted{\unsure{marqués}} \inserted{effectivement} marqués, \inserted{corr.} \uline{arêtes} sur $\Sigma$ \unsure{que}}
\margin{(\uline{NB} elles \unsure{sont les intersections|se trouvent sur|se déduisent} des cercles eff. marqués, et sont des intervalles ouverts)}
\mspage{55}
